%% Dans ce fichier, les �ventuels accents sont cod�s
%% en ISO-8859-15 (aka latin9)
%% La ligne suivante donne trois exemples 
%% e accent aigu : �, e accent grave : �, e-dans-l'o : �
%%%%%ENTREES%%%%
\begin{tabarg}
\multicolumn{3}{l}{\titretab {Entr\'ees}}\\
\hline
$X$ & Vecteur, \texttt{float} & Variable � interpoler ()\\
$Temps_X$ & Vecteur, \texttt{double} & Base de temps de $X$ (s)\\
$Temps_{Dest}$ & Vecteur, \texttt{double} & Base de temps d'arriv�e  (s)\\
$D$ & Coeff., \texttt{float} & D�calage � appliquer � $X$ (s)\\
$Unite_X$ & Coeff., \texttt{String} & Unit� de $X$ (vaut "radians", "degres" pour un angle) \\
$Nbval_{Dest}$ & Coeff., \texttt{Int\_32} & Nombre d'�chantillons dans la base de temps d'arriv�e  (dans fvb utiliser \$Arg3.NbVal)\\
\end{tabarg}

\medskip
%%%%%SORTIES%%%%
\begin{tabarg}
\multicolumn{3}{l}{\titretab {Sorties}}\\
\hline
$X_{sync}$ & Vecteur, \texttt{float} & Variable interpol�e et d�cal�e (m�me unit� que $X$)\\
\end{tabarg}
